\chapter{Introducción}
\label{cap:introduccion}

% -----------------------------------------------------------------------
% 1.1 CONTEXTO Y ANTECEDENTES
% -----------------------------------------------------------------------
\section{Contexto y antecedentes}
\label{sec:contexto}

La movilidad urbana constituye uno de los pilares fundamentales de la vida en las grandes ciudades modernas. En entornos metropolitanos como Madrid, millones de desplazamientos se producen cada día a pie, en bicicleta o en transporte público, convirtiendo la red viaria en un espacio compartido de enorme complejidad. Según datos del Ayuntamiento de Madrid, la ciudad cuenta con más de 3,3 millones de habitantes y una densidad de población que la sitúa entre las más elevadas de Europa, lo que intensifica la interacción entre peatones, ciclistas y vehículos motorizados a lo largo de sus más de 4.600 kilómetros de vías.

Históricamente, los sistemas de navegación y cálculo de rutas han priorizado criterios de eficiencia: minimizar la distancia recorrida o el tiempo de desplazamiento. Desde la introducción del GPS en 1973 \cite{bray2014} y la posterior proliferación de aplicaciones de navegación en teléfonos inteligentes, la planificación de rutas se ha consolidado como una tecnología ubicua. Sin embargo, tal y como señala la revisión sistemática de Sohrabi et al. (2022) \cite{sohrabi2022}, a pesar de los avances en los sistemas de navegación por carretera, la seguridad no está siendo considerada todavía en la planificación de rutas. Así, la ruta más rápida no es necesariamente la más segura, y en determinadas circunstancias puede incluso incrementar el riesgo de sufrir un accidente o un incidente de seguridad ciudadana.

La seguridad en el entorno urbano es un concepto multidimensional que abarca, al menos, tres grandes factores: la \textbf{seguridad vial} (riesgo de ser víctima de un atropello o colisión), la \textbf{seguridad ciudadana} (exposición a zonas con mayor índice de criminalidad) y la \textbf{seguridad ambiental} (condiciones del entorno que aumentan o reducen la percepción de riesgo, como la iluminación). En Madrid, el Ayuntamiento pone a disposición de la ciudadanía un conjunto de datos abiertos que permiten cuantificar la seguridad vial y la seguridad ambiental con una granularidad geográfica suficiente para integrarlos en un sistema de información geográfica (SIG). La seguridad ciudadana, en cambio, no puede cuantificarse a partir de estadísticas de criminalidad reales, ya que estas no se publican con desagregación geográfica por razones de protección de datos (véase la Sección~\ref{sec:motivacion}); este trabajo la aproxima mediante un índice de vulnerabilidad territorial, con las limitaciones que ello implica.

Los Sistemas de Información Geográfica representan, en este contexto, una herramienta de análisis especialmente adecuada. La combinación de bases de datos espaciales —como PostgreSQL con la extensión PostGIS— con motores de enrutamiento sobre grafos —como pgRouting— permite modelar la red viaria urbana como un grafo ponderado y calcular rutas óptimas bajo criterios de coste personalizados. Esta arquitectura tecnológica, ya consolidada en la comunidad científica y en proyectos de datos abiertos como OpenStreetMap, proporciona la infraestructura necesaria para desarrollar un sistema de enrutamiento que incorpore factores de seguridad como dimensión explícita del coste de cada arista del grafo.

% -----------------------------------------------------------------------
% 1.2 JUSTIFICACIÓN Y MOTIVACIÓN
% -----------------------------------------------------------------------
\section{Justificación y motivación}
\label{sec:motivacion}

La investigación de Sohrabi y Lord (2022), recogida en la revisión de Sohrabi et al. \cite{sohrabi2022}, demostró que la ruta más rápida no es necesariamente la más segura: un aumento del 8\% en el tiempo de desplazamiento puede asociarse a una reducción del 23\% en la probabilidad de verse involucrado en un accidente de tráfico. Este hallazgo pone de manifiesto la existencia de un \textit{trade-off} entre eficiencia temporal y seguridad que las aplicaciones de navegación convencionales ignoran sistemáticamente.

Desde la perspectiva de la seguridad ciudadana, trabajos como los de Galbrun et al. (2016) \cite{galbrun2016} o Levy et al. (2020) \cite{levy2020} han explorado la posibilidad de calcular rutas que minimicen la exposición a zonas con alta densidad de criminalidad, demostrando que es posible construir sistemas de enrutamiento seguros sin incurrir en desvíos excesivos respecto a la ruta más corta. En el ámbito peatonal específicamente, Lozano Domínguez y Mateo Sanguino (2021) \cite{lozano2021} propusieron un algoritmo de planificación de rutas seguras para peatones basado en lógica difusa, incorporando la detección de intención de cruce y otros factores de riesgo.

La motivación de este Trabajo Fin de Grado surge, por tanto, de la confluencia de tres elementos:

\begin{enumerate}
    \item \textbf{Una necesidad social real}: la seguridad en los desplazamientos urbanos a pie es una preocupación creciente, especialmente en determinados colectivos como mujeres que caminan solas de noche, personas mayores o turistas que no conocen la ciudad.
    \item \textbf{La disponibilidad de datos abiertos}: el Ayuntamiento de Madrid publica, a través de su portal de Datos Abiertos y del Geoportal, información georreferenciada sobre iluminación pública, estadísticas de accidentalidad y un índice de vulnerabilidad territorial por distrito, lo que hace técnicamente viable la construcción del sistema propuesto. No se ha encontrado, en cambio, una fuente pública de datos de criminalidad con desagregación geográfica: esta información la custodia el Ministerio del Interior y no se difunde a ese nivel de detalle por razones de protección de datos y de la propia investigación policial, lo que constituye una limitación asumida desde el planteamiento inicial del proyecto.
    \item \textbf{La madurez tecnológica del ecosistema SIG}: herramientas como PostGIS y pgRouting, junto con bibliotecas de visualización web como Leaflet, permiten construir prototipos funcionales de sistemas de enrutamiento seguro con tecnologías de código abierto y sin coste de licencia.
\end{enumerate}

En definitiva, este proyecto se justifica porque cubre un hueco real en el panorama de las aplicaciones de movilidad urbana: ofrece una alternativa a la ruta más corta o más rápida, poniendo en el centro la seguridad física del peatón mediante la integración de datos objetivos y georreferenciados.

% -----------------------------------------------------------------------
% 1.3 OBJETIVOS E HIPÓTESIS
% -----------------------------------------------------------------------
\section{Objetivos e hipótesis}
\label{sec:objetivos}

\subsection{Objetivo general}
\label{subsec:objetivo_general}

El objetivo general de este Trabajo Fin de Grado es \textbf{diseñar y desarrollar un sistema de cálculo de rutas seguras para peatones en el entorno urbano de Madrid}, que integre indicadores georreferenciados de seguridad (iluminación pública, vulnerabilidad territorial y accidentalidad) mediante técnicas de análisis espacial y enrutamiento sobre grafos ponderados, y que proporcione a los usuarios una alternativa informada a las rutas convencionales basadas únicamente en distancia o tiempo.

\subsection{Objetivos específicos}
\label{subsec:objetivos_especificos}

Para alcanzar el objetivo general, se plantean los siguientes objetivos específicos:

\begin{enumerate}
    \item \textbf{OE1 — Extracción y transformación de datos (ETL)}: recopilar, limpiar, normalizar y proyectar al mismo sistema de referencia de coordenadas los conjuntos de datos de iluminación pública, índice de vulnerabilidad territorial por distritos y accidentes de tráfico con víctimas peatonales proporcionados por el Ayuntamiento de Madrid.

    \item \textbf{OE2 — Modelado de la base de datos espacial}: diseñar e implementar un esquema de base de datos relacional-espacial en PostgreSQL/PostGIS que almacene la red viaria de Madrid (obtenida de OpenStreetMap), así como los indicadores de seguridad asociados a cada segmento de calle.

    \item \textbf{OE3 — Diseño de la función de coste de seguridad}: definir una función de ponderación que combine los indicadores de iluminación, vulnerabilidad territorial y accidentalidad en un coste de seguridad único asignable a cada arista del grafo viario, justificando las decisiones de diseño con fundamento bibliográfico.
    
    \item \textbf{OE4 — Implementación del algoritmo de enrutamiento seguro}: configurar pgRouting sobre la red ponderada y desarrollar las consultas SQL necesarias para calcular la ruta óptima en términos de seguridad entre dos puntos dados.
    
    \item \textbf{OE5 — Desarrollo de un prototipo de visor cartográfico}: construir una interfaz web ligera con Leaflet que permita al usuario introducir un origen y un destino, visualizar la ruta segura calculada sobre un mapa de Madrid y compararla con la ruta más corta convencional.
    
    \item \textbf{OE6 — Validación y análisis comparativo}: diseñar un conjunto de pruebas representativas que permitan evaluar cuantitativamente las diferencias entre la ruta segura y la ruta más corta en términos de longitud, tiempo estimado de recorrido y nivel de exposición al riesgo.
\end{enumerate}

\subsection{Hipótesis de partida}
\label{subsec:hipotesis}

La hipótesis principal que guía este proyecto es la siguiente:

\begin{quote}
\textit{Es posible calcular rutas peatonales que minimicen la exposición a factores de riesgo urbano (escasa iluminación, alta vulnerabilidad territorial y alta accidentalidad) con un incremento de distancia o tiempo de recorrido asumible respecto a la ruta más corta, mediante la integración de datos abiertos georreferenciados en un sistema de enrutamiento sobre grafos ponderados.}
\end{quote}

Esta hipótesis está apoyada por la evidencia existente en la literatura: Sohrabi et al. \cite{sohrabi2022} documentan que pequeños incrementos en el tiempo de viaje pueden asociarse a reducciones significativas en el riesgo, y trabajos como el de Mata et al. (2016) \cite{mata2016} han demostrado la viabilidad de sistemas similares en otras ciudades utilizando datos de crimen y accidentalidad.

% -----------------------------------------------------------------------
% 1.4 ESTRUCTURA DE LA MEMORIA
% -----------------------------------------------------------------------
\section{Estructura de la memoria}
\label{sec:estructura}

El presente documento se organiza en los siguientes capítulos:

\begin{itemize}
    \item \textbf{Capítulo 1 — Introducción} (presente capítulo): presenta el contexto del problema, la motivación del proyecto, los objetivos planteados y la organización de la memoria.
    
    \item \textbf{Capítulo 2 — Estado del arte}: revisa el dominio del problema desde el punto de vista científico, analizando los trabajos relacionados sobre enrutamiento seguro, los enfoques metodológicos existentes en la literatura y las tecnologías disponibles para su implementación.
    
    \item \textbf{Capítulo 3 — Descripción de la propuesta}: detalla la arquitectura del sistema propuesto, describiendo cada uno de sus componentes: el proceso ETL de datos, el modelado de la red viaria, la función de coste de seguridad, el motor de enrutamiento y el visor cartográfico.
    
    \item \textbf{Capítulo 4 — Metodología e implementación}: expone las decisiones técnicas adoptadas durante el desarrollo, incluyendo el diseño del esquema de base de datos espacial, las consultas de enrutamiento con pgRouting y el desarrollo del prototipo web.
    
    \item \textbf{Capítulo 5 — Gestión del proyecto}: recoge la planificación temporal del trabajo mediante un diagrama de Gantt, así como la estimación del presupuesto asociado al desarrollo.
    
    \item \textbf{Capítulo 6 — Pruebas y validación}: describe la metodología de evaluación utilizada y presenta los resultados del análisis comparativo entre la ruta segura y la ruta más corta.
    
    \item \textbf{Capítulo 7 — Conclusiones y trabajos futuros}: evalúa el grado de cumplimiento de los objetivos planteados, reflexiona sobre las limitaciones del sistema desarrollado y propone líneas de trabajo futuro.
    
    \item \textbf{Bibliografía}: recoge todas las referencias bibliográficas citadas a lo largo del documento.
    
    \item \textbf{Anexos}: incluye fragmentos de código relevantes, scripts ETL y consultas SQL/PostGIS de especial complejidad.
\end{itemize}
